\chapter{Conclusion générale et perspectives}
\label{chap:conclusion}

\section{Introduction}

Ce chapitre synthétise les apports du travail réalisé, en reliant explicitement les résultats techniques aux objectifs initiaux, au contexte de stage au sein de \#ADD, et aux exigences du cursus EMSI \cite{sommerville2016,bass2021}. Il propose ensuite des perspectives réalistes, classées par valeur ajoutée et complexité, afin d'éviter une liste de souhaits non fondée techniquement \cite{kleppmann2017,newman2021}.

\section{Bilan des objectifs}

Le travail réalisé couvre le cycle d'ingénierie d'une plateforme de services électroniques, depuis l'expression du besoin jusqu'à une proposition de déploiement reproductible \cite{pressman2014,ieee1012}. Les objectifs fonctionnels principaux — authentification, catalogue, cycle de vie de demande, rôles, notifications in-app — ont été traités de manière cohérente avec une architecture modulaire \cite{fowler2002,gamma1994}. Les objectifs non fonctionnels — sécurité de base, transactions, pagination, journalisation — ont guidé les choix d'implémentation et les tests associés \cite{owasp2021,iso25010}.

\section{Apports techniques}

Sur le plan technique, le projet démontre la maîtrise d'une stack moderne : Laravel pour une API structurée et une couche d'authentification SPA (Sanctum), React pour une interface interactive, PostgreSQL pour la persistance relationnelle, et Docker pour l'emballage lorsque l'environnement l'exige \cite{laraveldoc,sanctumdoc,reactdoc,postgresdoc,dockerdoc}. L'usage de migrations Artisan et de tests automatisés renforce la crédibilité de la démarche qualité \cite{kleppmann2017,sommerville2016}.

\section{Apports liés au stage chez \#ADD}

Le stage chez \#ADD a permis d'ancrer le projet dans des contraintes réelles : priorisation, validation par démonstration, collaboration avec l'équipe, et sensibilité à la maintenabilité \cite{bass2021}. L'encadrant de stage et l'équipe ont contribué à identifier tôt certains risques (transitions d'état, contrôle d'accès, gestion des fichiers) et à adopter des pratiques de revue utiles à la qualité finale \cite{pressman2014}.

\section{Limites reconnues}

Les limites incluent le périmètre volontairement réduit sur l'interopérabilité externe, l'absence de signature électronique qualifiée, et une politique de révocation JWT simplifiée pour démonstration \cite{rfc7519,owasp2021}. Les données réelles d'agence n'ont pas été exposées, ce qui limite certains indicateurs quantitatifs, mais respecte la confidentialité \cite{iso25010}.

\section{Leçons méthodologiques}

La leçon méthodologique majeure est l'intérêt d'une traçabilité continue entre besoins, conception et tests \cite{ieee1012,sommerville2016}. Lorsque cette traçabilité est tenue à jour, les arbitrages deviennent plus lisibles et la validation moins controversée \cite{pressman2014}. À l'inverse, lorsque les spécifications restent floues, le développement tend à improviser des règles métier implicites difficiles à maintenir \cite{fowler2002}.

\section{Perspectives fonctionnelles}

\subsection{Workflow avancé et paramétrage}

Une perspective majeure consiste à généraliser le moteur de workflow pour permettre des graphes d'états configurables par service, avec branches conditionnelles, SLA par étape, et escalades automatiques \cite{bass2021}. Cette évolution nécessite un meta-modèle plus riche et une interface d'administration dédiée \cite{sommerville2016}.

\subsection{Interopérabilité et intégrations}

L'intégration avec des annuaires d'identité fédérés, des bus d'événements, ou des API tierces permettrait de réduire la saisie et d'améliorer la fiabilité des référentiels \cite{kleppmann2017,newman2021}. L'interopérabilité impose une gouvernance contractuelle stricte (schémas, versionnement, résilience) \cite{fielding2000}.

\section{Perspectives techniques et opérationnelles}

\subsection{Observabilité et exploitation}

L'ajout de métriques, de traces distribuées, et de tableaux de bord améliorerait l'exploitabilité \cite{bass2021}. La corrélation logs/traces faciliterait le diagnostic des incidents en production \cite{iso25010}.

\subsection{Performance et scalabilité}

Au-delà des tests légers, une démarche de performance engineering pourrait inclure profilage JDBC, optimisation de requêtes, mise en cache sélective, et éventuellement lecture répliquée pour des usages analytiques \cite{kleppmann2017,postgresdoc}.

\subsection{Industrialisation CI/CD}

Une chaîne CI/CD complète pourrait intégrer : build parallèle, tests, analyse statique, publication d'images, déploiement sur environnement cloud, et stratégie blue/green ou canary selon criticité \cite{dockerdoc,twelvefactor,sommerville2016}.

\section{Perspectives sécurité et conformité}

Les perspectives sécurité incluent rotation automatique des clés, stockage des secrets dans un coffre, durcissement des en-têtes HTTP, politique CSP stricte côté frontend, scans de dépendances en CI, et modélisation explicite des menaces \cite{owasp2021}. Sur le plan conformité, une analyse RGPD détaillée devrait accompagner toute mise en production traitant des données personnelles \cite{iso25010}.

\section{Perspectives UX et accessibilité}

Des optimisations de performance perçue (pagination virtuelle, chargement paresseux) et un audit d'accessibilité (WCAG) renforceraient l'inclusivité et la qualité perçue \cite{pressman2014,iso25010}.

\section{Ouverture recherche et innovation}

Des pistes de recherche incluent l'usage d'assistants décisionnels pour trier certaines demandes simples, sous réserve d'encadrement humain et de conformité réglementaire \cite{sommerville2016}. Une telle perspective dépasse le MVP et nécessite une gouvernance éthique rigoureuse \cite{iso25010}.

\section{Conclusion finale}

En synthèse, le mémoire a proposé une solution complète sur le plan documentaire et technique, ancrée dans des technologies éprouvées, avec une attention particulière portée à la sécurité des accès et à l'intégrité des transitions d'état \cite{owasp2021,date2003}. Le dispositif constitue une base solide pour des extensions fonctionnelles et opérationnelles, en prolongeant les apprentissages du stage au sein de \#ADD et les exigences de formation de l'EMSI Rabat \cite{bass2021,sommerville2016}.
